\beginsong{Love}[by={Jean-Claude Gianadda}]
\textnote{Capo 5}
\beginverse
\[Am]Il est venu chez nous, gui\[Dm]tare en bandoulière,
Venait d'on ne sait \[G]où, il parcourait la \[C]terre.
Et dans ses longs che\[E]veux, le vent semblait chan\[Am]ter,
Tout au fond de ses \[E]yeux dansait la liber\[Am]té.
\endverse
\beginchorus
\[Am]Love, c'était son \[Dm]nom, \[G]Love,
Un vaga\[C]bond qui vivait de so\[E]leil,
D'espace et de chan\[Am]sons.
\endchorus
\beginverse
\[Am]Il écoutait le vent, les \[Dm]fleurs et les rivières,
Jouait comme un en\[G]fant, parlait de la lu\[C]mière.
Il partageait ses \[E]rires, ses rêves et ses pro\[Am]jets,
Et dans chaque sou\[E]rire dansait la liber\[Am]té.
\endverse
\beginverse
\[Am]Il est parti un jour, \[Dm]nul ne sait où il est,
Au pays de l'a\[G]mour, tu peux le rencon\[C]trer.
Et dans notre mai\[E]son, il nous aura lais\[Am]sé,
Avec quelques chan\[E]sons, un peu de liber\[Am]té.
\endverse
\endsong
